\section{Introduction}
\label{sec:intro}

Design and analysis loops---optimization, uncertainty propagation, calibration, reliability---call a simulation many
thousands of times. When one call costs minutes to hours, as for the coupled finite-element analyses behind the present
work, the loop is run on a \emph{surrogate}: a cheap model fitted to a limited set of simulation results
\cite{queipo2005surrogate,forrester2008engineering,forrester2009recent}. The idea has been reinvented under many names
in many fields---response surfaces in chemical engineering \cite{box1951experimental}, kriging in mining geostatistics
\cite{krige1951statistical,matheron1963principles,cressie1990origins}, design and analysis of computer experiments in
statistics \cite{sacks1989design,santner2003design}, proxy models in petroleum engineering \cite{bahrami2022proxy},
emulators in climate science, cosmology and epidemiology
\cite{watsonparris2022climatebench,heitmann2009coyote,andrianakis2015bayesian}, metamodels in simulation and operations
research \cite{kleijnen2015design}, surrogate benchmarks in automated machine learning
\cite{eggensperger2015efficient}, and fitness approximation in evolutionary computation \cite{jin2005comprehensive,jin2011surrogate}.

Each community has converged on its own default. In engineering design the Gaussian process (GP) is the dominant
surrogate \cite{forrester2008engineering,gramacy2020surrogates}, and multi-fidelity GP formulations are the natural
extension when cheaper analyses are available \cite{kennedy2000predicting,forrester2007multi,needels2024efficient}. In
tabular machine learning, tree ensembles remain state of the art at small-to-medium sample sizes
\cite{grinsztajn2022tree,shwartzziv2022tabular,mcelfresh2023neural}, and tree-based surrogates are preferred for
algorithm-configuration benchmarks \cite{eggensperger2015efficient,eggensperger2018efficient} and agent-based model
calibration \cite{lamperti2018agent}. These defaults are rarely confronted with each other on the same problems, with
the same budgets, under the same noise.

The confrontation is not academic. The surrogate study of Pozna\'nski \cite[Sec.~14.3 and App.~G]{poznanski2026towards},
which sized the deployable duct wall of a tube-launched ramjet from coupled fluid--structure finite-element analyses,
reported a GP with the lowest error on six synthetic surfaces, yet gradient-boosted trees with the lowest
cross-validated error on the real finite-element data; the trees were set aside because the downstream interior-point
optimizer requires smooth derivatives. That study selected its surrogate by approximation accuracy, implicitly assuming
that the most accurate surrogate yields the best design---an assumption contradicted by several benchmarks in which the
approximation winner and the optimization winner differ
\cite{williams2021selection,eggensperger2015efficient,neufang2024surrogate}. The present paper revisits and extends that
study; where its findings are corrected here, the corrections are stated explicitly.

This paper has two parts. The \textbf{review} (Sections~\ref{sec:landscape}--\ref{sec:theory}) surveys, breadth first and
across disciplines, the methods that make up a surrogate pipeline: the learner (Section~\ref{sec:landscape}), its
uncertainty estimate and the sampling it drives (Section~\ref{sec:uq}), multi-fidelity fusion (Section~\ref{sec:mf}),
change-of-basis and latent-space representations borrowed from data-driven reduced-order modelling
(Section~\ref{sec:latent}), the conversion of a learner into an optimizer-ready representation
(Section~\ref{sec:upscaling}), and surrogate-based optimization itself (Section~\ref{sec:sbo}), followed by the theory
that predicts when each should win (Section~\ref{sec:theory}). We deliberately include old methods that were not pursued
and methods that are standard outside aerospace but rare within it. The \textbf{study} (in preparation) varies surface
class, dimension, budget, noise level and type, fidelity structure and cost, scout--learner pairing, upscaling and
optimizer type in a factorial design, and scores pipelines by true-function regret per unit simulation cost.
Section~\ref{sec:gap} states the gap the study fills.
